% =====================================================================
% Experimentación y resultados — Agentes especializados y Agent Skills
% Datos: Experimentos/experiments/experimento_21_agentskills/results/
% Corpus de evaluación: 1.313 noticias anotadas en las cinco variables
% =====================================================================

La pregunta que vertebra toda la experimentación es una: \textbf{¿aporta algo
empaquetar la metodología experta como \textit{Agent Skills}?} Para responderla se
enfrenta la arquitectura propuesta (\textbf{nivel B1}) a una condición de control
que recibe la misma metodología por la vía más simple, inyectada íntegramente en el
\textit{prompt} (\textbf{nivel B0}). Ambos niveles se definen en el
Capítulo~\ref{implementation} (Sección~\ref{subsec:diseno_agentico}).

Aislar el efecto de las \textit{skills} exige que todo lo demás permanezca fijo.
Por eso B0 y B1 comparten el mismo corpus, los mismos
modelos, las mismas cinco variables y el mismo formato de salida: la única
diferencia entre ambos es cómo llega la metodología al modelo, inyectada de una vez
en el \textit{prompt} (B0) o cargada bajo demanda por el agente (B1), de modo que
cualquier variación en el acuerdo es atribuible al empaquetado como \textit{Agent
Skills} y no a otro factor. Para homogeneizar además el comportamiento entre
proveedores, se desactivó el razonamiento explícito en todos los modelos
(\texttt{reasoning\_effort} en OpenAI, \texttt{thinking\_budget} en Google). La
interpretación de las métricas exige tener presente la prevalencia de cada variable
y la exactitud trivial asociada, recogidas en la
Tabla~\ref{tab:iris-prevalencia-imp} del Capítulo~\ref{sec:methodology}.

\subsection{Resultados del núcleo: \textit{baseline} frente a Agent Skills}
\label{subsec:iris_core}

La Tabla~\ref{tab:iris-b0b1} compara ambos niveles en los cuatro modelos
evaluados: los tres accesibles por API de pago y el modelo local
\texttt{gemma4:e4b}.

\begin{table}[H]
    \centering
    \begingroup
    \scriptsize
    \setlength{\tabcolsep}{4pt}
    \newcommand{\dr}[1]{\textcolor{gray!75}{$#1$}} % fila delta: en gris
    \renewcommand{\texttt}[1]{{\ttfamily #1}}
    \ttabbox[\FBwidth]{
        \caption{Comparación B0 (\textit{baseline}) frente a B1 (Agent Skills). Medias sobre las cinco variables, $N=1.313$. Las filas $\Delta$ (en gris) recogen la variación B1$-$B0.}
        \label{tab:iris-b0b1}
    }{
        \begin{tabularx}{\textwidth}{
            >{\hsize=1.7\hsize\raggedright\arraybackslash}X
            >{\hsize=0.6\hsize\centering\arraybackslash}X
            >{\hsize=0.925\hsize\centering\arraybackslash}X
            >{\hsize=0.925\hsize\columncolor{gray!12}\centering\arraybackslash}X
            >{\hsize=0.925\hsize\centering\arraybackslash}X
            >{\hsize=0.925\hsize\centering\arraybackslash}X
        }
            \toprule
            \textbf{Modelo} & \textbf{Nivel} & \textbf{Exactitud} & \boldmath$\kappa$ & \textbf{F1 macro} & \textbf{Coste (USD)} \\
            \midrule
            \multirow{3}{*}{\texttt{gemini-3.1-flash-lite}}
              & B0 & 0,600 & 0,026 & 0,388 & 5,64 \\
              & B1 & 0,661 & \textbf{0,066} & 0,443 & 13,16 \\
              & \dr{\Delta} & \dr{+0,061} & \dr{+0,040} & \dr{+0,055} & \dr{2{,}3\times} \\
            \midrule
            \multirow{3}{*}{\texttt{gpt-4o-mini}}
              & B0 & 0,564 & 0,033 & 0,410 & 2,52 \\
              & B1 & 0,612 & 0,014 & 0,438 & 6,89 \\
              & \dr{\Delta} & \dr{+0,048} & \dr{-0,019} & \dr{+0,029} & \dr{2{,}7\times} \\
            \midrule
            \multirow{3}{*}{\texttt{gpt-5.4-nano}}
              & B0 & 0,571 & 0,052 & 0,430 & 5,67 \\
              & B1 & 0,569 & 0,004 & 0,364 & 10,86 \\
              & \dr{\Delta} & \dr{-0,003} & \dr{-0,048} & \dr{-0,066} & \dr{1{,}9\times} \\
            \midrule
            \multirow{2}{*}{\textbf{Media (API)}}
              & \textbf{B0} & 0,578 & \textbf{0,037} & 0,409 & \textemdash \\
              & \textbf{B1} & \textbf{0,614} & 0,028 & 0,415 & \textemdash \\
            \midrule
            \multirow{3}{*}{\texttt{gemma4:e4b}\,\textnormal{\footnotesize(local)}}
              & B0 & 0,597 & \textbf{0,061} & 0,421 & \textemdash \\
              & B1 & 0,605 & 0,043 & 0,412 & \textemdash \\
              & \dr{\Delta} & \dr{+0,008} & \dr{-0,018} & \dr{-0,009} & \textemdash \\
            \bottomrule
        \end{tabularx}
    }
    \endgroup
\end{table}

A primera vista, la arquitectura parece funcionar: la exactitud sube con B1 en la
mayoría de los modelos (Figura~\ref{fig:iris-b0b1-metricas}), y la exactitud media de
las API pasa de 0,578 a 0,614. Pero exactitud y $\kappa$ \emph{no cuentan la
misma historia}: el $\kappa$ medio, lejos de acompañar, baja de 0,037 a 0,028, y todo
ello a un coste de inferencia 2,2 veces mayor. Esta discordancia obliga a precisar
qué métrica resulta pertinente antes de interpretar el efecto de las \textit{skills}:
en una tarea con clases fuertemente desbalanceadas, la exactitud es un indicador
engañoso y el análisis debe apoyarse en el coeficiente $\kappa$, según se argumenta
en la Sección~\ref{subsec:iris_divergencia}.

\begin{figure}[H]
    \centering
    \includegraphics[width=\textwidth]{fig_b0b1_barras_blanco.pdf}
    \caption{Niveles de exactitud, F1 macro y $\kappa$ por modelo, en B0 y B1
    ($N=1.313$). Cada métrica se representa en su propia escala: el $\kappa$ vive
    en un orden de magnitud mucho menor que la exactitud.}
    \label{fig:iris-b0b1-metricas}
\end{figure}


\subsection{Divergencia entre exactitud y \texorpdfstring{$\kappa$}{kappa}}
\label{subsec:iris_divergencia}

Conviene detenerse en esa discordancia, porque condiciona cómo debe leerse toda la
tabla: \textbf{la mejora en exactitud no se traduce en más acuerdo}. Considerando
las medias, el nivel B1 sube la exactitud en 0,036 pero el $\kappa$ baja en $-0,009$.
Exactitud y $\kappa$ pueden, por tanto, moverse de forma independiente e incluso en
sentidos opuestos. Esto es un efecto conocido de $\kappa$ en
presencia de clases desbalanceadas~\cite{feinstein1990}, pero aquí resulta determinante: pasarlo por alto
invertiría la conclusión del capítulo. La Figura~\ref{fig:iris-divergencia} lo
visualiza como trayectorias B0$\to$B1: en el panel de exactitud casi todas las líneas
ascienden, mientras que en el de $\kappa$ se mantienen o descienden, salvo en
\texttt{gemini}, cuya línea sube en ambas.

\begin{figure}[H]
    \centering
    \includegraphics[width=\textwidth]{fig_b0b1_slope_blanco.pdf}
    \caption{Evolución de cada modelo de B0 a B1 en las tres métricas ($N=1.313$).
    La exactitud sube en tres de los cuatro modelos, pero el $\kappa$ solo mejora en
    \texttt{gemini-3.1-flash-lite}. En el resto se mantiene o desciende. La mejora en
    exactitud sin mejora en $\kappa$ es el desacoplamiento entre exactitud y acuerdo.}
    \label{fig:iris-divergencia}
\end{figure}

El motivo de esa divergencia es conocido. Al volverse menos conservador, el modelo
responde \textit{Sí} con más frecuencia, y cuando la clase positiva es la mayoritaria,
contestar afirmativamente más a menudo eleva la exactitud por simple coincidencia
con esa mayoría, no por discriminar mejor. El coeficiente $\kappa$, que descuenta el
acuerdo esperable por azar, no premia ese desplazamiento. La consecuencia práctica
es directa: \emph{un estudio que reportara solo exactitud concluiría que las Agent
Skills mejoran el sistema, y se equivocaría}. De ahí la recomendación de reportar
siempre exactitud, $\kappa$ y prevalencia de forma conjunta en tareas de
clasificación desbalanceadas.

Leído correctamente, por $\kappa$, el efecto de las \textit{Agent Skills}
\textbf{depende del modelo}. Solo \texttt{gemini-3.1-flash-lite} mejora con la
arquitectura: su $\kappa$ pasa de 0,026 a 0,066, y lo hace acompañado de subidas
simultáneas en exactitud y F1. En los dos modelos de OpenAI las \textit{skills} no ayudan e incluso
perjudican ($\Delta\kappa = -0,019$ en \texttt{gpt-4o-mini} y $-0,048$ en
\texttt{gpt-5.4-nano}), y el modelo local se comporta igual que ellos ($\kappa$ de
0,061 a 0,043). De los cuatro modelos, solo uno mejora con las \textit{skills}. En
los otros tres el $\kappa$ se mantiene o desciende, y la media no aumenta. Que el
efecto dependa del modelo y no eleve el acuerdo en agregado admite varias lecturas, desde el techo de la propia tarea hasta la heterogeneidad de la anotación, que se
retoman en la discusión (Sección~\ref{sec:discussion}) una vez presentada toda la
evidencia.

\subsection{Descomposición de la arquitectura por componentes}
\label{subsec:iris_descomposicion}

El nivel B1 combina tres ingredientes, descritos en el Capítulo~\ref{implementation}:
la carga de \textit{skills} bajo demanda con verificación de evidencias
(Sección~\ref{subsec:diseno_agentico}), un catálogo de resúmenes de las guías
expertas descubrible por el agente (Sección~\ref{prompts}), y una herramienta de
recuperación en vivo sobre esas guías (RAG, implementada como \texttt{CONSULTAR\_GUIA},
Sección~\ref{rag}). Para averiguar cuál de ellos explica el efecto
observado se ejecutaron cinco variantes por modelo, activando y desactivando cada
componente de forma independiente, sobre los cuatro modelos y las 1.313 noticias.
Las Tablas~\ref{tab:iris-descomposicion}, \ref{tab:iris-descomposicion-acc}
y~\ref{tab:iris-descomposicion-f1} recogen, respectivamente, el $\kappa$, la
exactitud y el F1 macro medios de cada combinación.

\begin{table}[H]
    \centering
    \begingroup
    \scriptsize
    \ttabbox[\FBwidth]{
        \caption{Descomposición de B1 en componentes. $\kappa$ medio sobre las cinco variables, por modelo ($N=1.313$)}
        \label{tab:iris-descomposicion}
    }{
        \begin{tabularx}{\textwidth}{
            X
            >{\centering\arraybackslash}X
            >{\centering\arraybackslash}X
            >{\centering\arraybackslash}X
            >{\centering\arraybackslash}X
        }
            \toprule
            \textbf{Configuración} & \textbf{Gemini} & \textbf{GPT-4o-mini} & \textbf{GPT-5.4-nano} & \textbf{Gemma} \\
            \midrule
            B0 & 0,026 & 0,033 & 0,052 & 0,061 \\
            Solo \textit{skills} & 0,061 & 0,025 & 0,000 & 0,065 \\
            \textit{Skills} + resúmenes & \textbf{0,075} & 0,024 & 0,015 & 0,061 \\
            \textit{Skills} + RAG & 0,063 & 0,028 & 0,012 & 0,038 \\
            Completo & 0,066 & 0,014 & 0,004 & 0,043 \\
            \bottomrule
        \end{tabularx}
    }
    \endgroup
\end{table}

\begin{table}[H]
    \centering
    \begingroup
    \scriptsize
    \ttabbox[\FBwidth]{
        \caption{Descomposición de B1 en componentes. \textbf{Exactitud} media sobre las cinco variables, por modelo ($N=1.313$)}
        \label{tab:iris-descomposicion-acc}
    }{
        \begin{tabularx}{\textwidth}{
            X
            >{\centering\arraybackslash}X
            >{\centering\arraybackslash}X
            >{\centering\arraybackslash}X
            >{\centering\arraybackslash}X
        }
            \toprule
            \textbf{Configuración} & \textbf{Gemini} & \textbf{GPT-4o-mini} & \textbf{GPT-5.4-nano} & \textbf{Gemma} \\
            \midrule
            B0 & 0,600 & 0,564 & 0,571 & 0,597 \\
            Solo \textit{skills} & 0,645 & \textbf{0,677} & 0,567 & 0,608 \\
            \textit{Skills} + resúmenes & 0,656 & 0,609 & 0,569 & 0,610 \\
            \textit{Skills} + RAG & 0,651 & 0,631 & 0,571 & 0,602 \\
            Completo & 0,661 & 0,612 & 0,569 & 0,605 \\
            \bottomrule
        \end{tabularx}
    }
    \endgroup
\end{table}

\begin{table}[H]
    \centering
    \begingroup
    \scriptsize
    \ttabbox[\FBwidth]{
        \caption{Descomposición de B1 en componentes. \textbf{F1 macro} media sobre las cinco variables, por modelo ($N=1.313$)}
        \label{tab:iris-descomposicion-f1}
    }{
        \begin{tabularx}{\textwidth}{
            X
            >{\centering\arraybackslash}X
            >{\centering\arraybackslash}X
            >{\centering\arraybackslash}X
            >{\centering\arraybackslash}X
        }
            \toprule
            \textbf{Configuración} & \textbf{Gemini} & \textbf{GPT-4o-mini} & \textbf{GPT-5.4-nano} & \textbf{Gemma} \\
            \midrule
            B0 & 0,388 & 0,410 & 0,430 & 0,421 \\
            Solo \textit{skills} & 0,430 & \textbf{0,484} & 0,364 & 0,424 \\
            \textit{Skills} + resúmenes & 0,444 & 0,451 & 0,365 & 0,423 \\
            \textit{Skills} + RAG & 0,439 & 0,473 & 0,374 & 0,409 \\
            Completo & 0,443 & 0,438 & 0,364 & 0,412 \\
            \bottomrule
        \end{tabularx}
    }
    \endgroup
\end{table}

\begin{figure}[H]
    \centering
    \includegraphics[width=\textwidth]{fig_kappa_componentes_blanco.pdf}
    \caption{Coeficiente $\kappa$ medio al activar cada componente de B1, por modelo
    ($N=1.313$). La mejor configuración de cada modelo prescinde de la recuperación
    en vivo (RAG, franja sombreada). El mecanismo de \textit{skills} solo eleva el
    acuerdo en \texttt{gemini-3.1-flash-lite}.}
    \label{fig:iris-descomposicion-fig}
\end{figure}

Dos conclusiones se sostienen en los cuatro modelos, y la
Figura~\ref{fig:iris-descomposicion-fig} las resume. La primera afecta a la
herramienta RAG de recuperación en vivo: no aporta, y con frecuencia perjudica. En
todos los modelos la mejor variante prescinde de ella. En \texttt{gemma4:e4b} el
contraste es nítido: sin la RAG el $\kappa$ se sitúa entre 0,061 y
0,065, y al activarla cae a valores entre 0,038 y 0,043. La recuperación
por TF-IDF sobre textos periodísticos cortos en este caso introduce más ruido que señal.

La segunda conclusión concierne al mecanismo de \textit{skills}: solo mejora el
acuerdo en \texttt{gemini}. En ese modelo, la variante que se limita a cargar la
metodología como \textit{skill} verificable, sin guías de ningún tipo, ya recupera
la mayor parte de la mejora sobre el control, de 0,026 a 0,061. En los otros
tres modelos esa misma variante iguala el control (\texttt{gemma}) o lo empeora
(los dos de OpenAI). Cuando hay ganancia, la configuración con mejor acuerdo es
\textit{skills} más resúmenes sin RAG: más simple y más barata que el
sistema completo, y en ningún modelo peor.

En conjunto, la descomposición confirma la lectura del núcleo: la arquitectura no es
útil ni inútil en abstracto, sino que \textbf{amplifica lo que cada modelo ya trae}.
Aporta cuando el modelo aprovecha de verdad la carga bajo demanda, y estorba cuando la
infrautiliza o se satura con el catálogo. Su valor está condicionado a la capacidad
del modelo para explotarla.

El desglose en exactitud y F1 macro matiza el cuadro y, de paso, ilustra la
divergencia de la Sección~\ref{subsec:iris_divergencia}: la exactitud puede subir con
las \textit{skills} (\texttt{gpt-4o-mini} con solo \textit{skills} pasa de 0,564 a
0,677) sin que el acuerdo la acompañe (su $\kappa$ se queda en 0,025).


\subsection{Diagnóstico por variable}
\label{subsec:iris_diagnostico}

Hasta ahora el análisis se ha apoyado en métricas globales, promediadas sobre las
cinco variables. Ese nivel agregado responde a la pregunta central del capítulo,
pero esconde comportamientos muy distintos de una variable a otra. Para entender qué
ocurre por debajo del promedio hay que desglosar los resultados variable a variable.
Para cada una se reportan, sobre la clase positiva (\textit{Sí}), la \textbf{precisión} (de
las piezas que el modelo marca como \textit{Sí}, cuántas lo son de verdad, que no debe
confundirse con la exactitud, la cual mide el acierto global sobre todas las piezas)
y el \textbf{recall} o sensibilidad (del sexismo que marcan las expertas, cuánto
detecta el modelo), junto al coeficiente $\kappa$. El recall es aquí el más revelador: un recall bajo significa
que el modelo deja sin marcar buena parte de lo que las anotadoras sí marcan. La
Tabla~\ref{tab:iris-diagnostico} recoge estos valores para el nivel B1, combinándolos con la
revisión manual de casos se caracteriza cada variable.

\begin{table}[H]
    \centering
    \begingroup
    \scriptsize
    \setlength{\tabcolsep}{3pt}
    \newcommand{\tri}[3]{\begin{tabular}[t]{@{}c@{}}#1\\#2\\#3\end{tabular}}
    \ttabbox[\FBwidth]{
        \caption{Precisión, recall (sensibilidad) sobre la clase \textit{Sí} y $\kappa$, por variable en el nivel B1 ($N=1.313$). En cada celda se apilan, de arriba abajo, precisión, recall y $\kappa$. ``Prev.'' es la prevalencia de la clase \textit{Sí}.}
        \label{tab:iris-diagnostico}
    }{
        \begin{tabularx}{\textwidth}{
            >{\hsize=1.35\hsize\raggedright\arraybackslash\ttfamily}X
            >{\hsize=0.55\hsize\centering\arraybackslash}X
            >{\hsize=0.775\hsize\centering\arraybackslash}X
            >{\hsize=0.775\hsize\centering\arraybackslash}X
            >{\hsize=0.775\hsize\centering\arraybackslash}X
            >{\hsize=0.775\hsize\centering\arraybackslash}X
        }
            \toprule
            \normalfont\textbf{Variable} & \textbf{Prev.} & \textbf{Gemini} & \textbf{GPT-4o-mini} & \textbf{GPT-5.4-nano} & \textbf{Gemma} \\
            \midrule
            lenguaje\_sexista & 0,810 & \tri{0,961}{0,046}{0,015} & \tri{0,789}{0,450}{$-$0,037} & \tri{1,000}{0,001}{0,000} & \tri{0,929}{0,012}{0,003} \\
            \midrule
            masc\_generico & 0,810 & \textbf{\tri{0,887}{0,733}{0,261}} & \tri{0,873}{0,206}{0,034} & \tri{0,874}{0,104}{0,016} & \tri{0,897}{0,408}{0,109} \\
            \midrule
            sexismo\_discurso & 0,428 & \tri{0,423}{0,020}{0,000} & \tri{0,450}{0,112}{0,011} & \tri{0,435}{0,114}{0,004} & \tri{0,378}{0,025}{$-$0,006} \\
            \midrule
            asimetria\_mujer\_hombre & 0,062 & \tri{0,118}{0,074}{0,045} & \tri{0,063}{0,148}{0,003} & \tri{0,000}{0,000}{$-$0,012} & \tri{0,143}{0,160}{0,092} \\
            \midrule
            denominacion\_sexualizada & 0,100 & \tri{0,167}{0,015}{0,011} & \tri{0,375}{0,046}{0,061} & \tri{1,000}{0,008}{0,014} & \tri{0,158}{0,023}{0,015} \\
            \bottomrule
        \end{tabularx}
    }
    \endgroup
\end{table}

\begin{itemize}
    \item \textbf{\texttt{lenguaje\_sexista} (V25): un criterio en disputa.} La
    precisión es alta cuando el modelo marca \textit{Sí} (entre 0,79 y 1,00), pero el
    $\kappa$ ronda el azar en los cuatro modelos (entre $-0,037$ y 0,015). El problema
    no es que el modelo se equivoque al marcar, sino que sistema y anotadoras no
    coinciden en \emph{cuándo} hay que marcar. La causa está en el propio criterio de
    V25, que es intrínsecamente amplio y admite dos lecturas. Al ser la variable más
    general del bloque (el uso discriminatorio del lenguaje por razón de sexo, sin
    acotar a un mecanismo concreto), las anotadoras se inclinan por una lectura
    inclusiva y la marcan en casi todas las piezas, de ahí su prevalencia del
    81\,\%, mientras que el modelo se atiene a una lectura más literal y la marca
    mucho menos, de ahí su recall bajísimo. Ninguna de las dos lecturas es
    incorrecta, porque el propio criterio admite ambas.

    \item \textbf{\texttt{masc\_generico} (V26): el único acuerdo genuino, y su
    umbral.} Es el único caso con acuerdo apreciable: \texttt{gemini} alcanza
    $\kappa = 0,261$ y un recall de 0,733, es decir, detecta la mayoría del masculino
    genérico que marcan las expertas, con exactitud y $\kappa$ que mejoran a la vez
    sobre el \textit{baseline} (lo que descarta el artefacto de marginales y confirma
    una ganancia real). El local \texttt{gemma4:e4b} queda en segundo lugar, con un
    recall de 0,408 y un $\kappa$ de 0,109, la única variable en la que también supera
    con claridad el azar. Los dos modelos de OpenAI, en cambio, tienen un recall mucho
    menor (0,206 y 0,104): reflejan un desajuste de umbral, porque el español emplea el
    masculino genérico de forma pervasiva (\textit{los trabajadores}, \textit{los
    ciudadanos}) y la anotación tiende a marcarlo de forma sistemática, mientras que
    los modelos solo lo señalan cuando resulta evitable e invisibilizador.

    \item \textbf{\texttt{sexismo\_discurso} (V30): anotación sobre-inclusiva.} El
    acuerdo ronda el azar en los cuatro modelos ($\kappa$ entre $-0,006$ y 0,011) y el
    recall es muy bajo (entre 0,020 y 0,114). En una fracción de los casos positivos
    no se aprecia, en una lectura atenta del texto, el fenómeno codificado. Estos
    casos se concentran en el anotador de menor acuerdo global (véase la
    Sección~\ref{subsec:iris_errores}).

    \item \textbf{\texttt{asimetria\_mujer\_hombre} (V33) y
    \texttt{denominacion\_sexualizada} (V35): desbalance extremo.} Con prevalencias
    del 6,2\,\% y el 10,0\,\%, el modelo casi no marca positivos (recall entre 0,000 y
    0,160) y la precisión, cuando no es nula, descansa sobre un puñado de
    predicciones: \texttt{gpt-5.4-nano} marca \textit{Sí} en \texttt{denominacion\_sexualizada}
    tan pocas veces que su precisión de 1,000 carece de valor. En este régimen
    $\kappa$ penaliza severamente incluso a clasificadores razonables, hasta el punto
    de que una exactitud de 0,90 puede coexistir con $\kappa = 0,000$.
\end{itemize}

El desglose por variable revela además dónde actúan las \textit{skills}. La
Tabla~\ref{tab:iris-porvariable-b0b1} compara el $\kappa$ de B0 y B1 en cada
variable. En \texttt{gemini}, la mejora del núcleo se concentra por completo en
\texttt{masc\_generico} (de 0,073 a 0,261), mientras el resto de variables apenas se
mueve. En los otros tres modelos ocurre lo contrario: las \textit{skills} degradan
variables donde el control ya tenía algo de señal, sobre todo
\texttt{denominacion\_sexualizada} (de 0,123 a 0,014 en \texttt{gpt-5.4-nano}, de
0,136 a 0,061 en \texttt{gpt-4o-mini} y de 0,154 a 0,015 en \texttt{gemma4:e4b}). El
efecto de la arquitectura no solo depende del modelo, sino que dentro de cada modelo
se localiza en variables concretas.

\begin{table}[H]
    \centering
    \begingroup
    \scriptsize
    \ttabbox[\FBwidth]{
        \caption{Coeficiente $\kappa$ por variable en B0 (\textit{baseline}) y B1 (Agent Skills), para los cuatro modelos ($N=1.313$)}
        \label{tab:iris-porvariable-b0b1}
    }{
        \begin{tabularx}{\textwidth}{
            >{\hsize=1.4\hsize\raggedright\arraybackslash\ttfamily}X
            >{\hsize=0.9\hsize\centering\arraybackslash}X
            >{\hsize=0.9\hsize\centering\arraybackslash}X
            >{\hsize=0.9\hsize\centering\arraybackslash}X
            >{\hsize=0.9\hsize\centering\arraybackslash}X
        }
            \toprule
            \normalfont\textbf{Variable} & \textbf{Gemini} & \textbf{GPT-4o-mini} & \textbf{GPT-5.4-nano} & \textbf{Gemma} \\
             & \multicolumn{4}{c}{\normalfont\scriptsize $\kappa$: B0 / B1} \\
            \midrule
            lenguaje\_sexista & 0,000 / 0,015 & $-0,013$ / $-0,037$ & 0,001 / 0,000 & $-0,004$ / 0,003 \\
            masc\_generico & \textbf{0,073 / 0,261} & 0,023 / 0,034 & 0,103 / 0,016 & 0,102 / 0,109 \\
            sexismo\_discurso & 0,004 / 0,000 & 0,002 / 0,011 & 0,007 / 0,004 & 0,002 / $-0,006$ \\
            asimetria\_mujer\_hombre & 0,036 / 0,045 & 0,019 / 0,003 & 0,026 / $-0,012$ & 0,054 / 0,092 \\
            denominacion\_sexualizada & 0,019 / 0,011 & 0,136 / 0,061 & 0,123 / 0,014 & 0,154 / 0,015 \\
            \bottomrule
        \end{tabularx}
    }
    \endgroup
\end{table}

Conviene mirar también cuánto marca cada modelo en términos absolutos, que es el
reverso del recall: marcar poco implica dejar sin detectar buena parte de lo que las
expertas sí señalan. La Tabla~\ref{tab:iris-prev-modelo-var} recoge, por variable, la
prevalencia de \textit{Sí} de cada modelo. La única variable en la que algún modelo
se acerca al terreno de las expertas es \texttt{masc\_generico} (\texttt{gemini},
66,9\,\%), mientras que en el resto todos marcan \textit{Sí} en una fracción muy
pequeña de las piezas. El patrón, además, no es común: \texttt{gpt-4o-mini} dispara
en \texttt{lenguaje\_sexista} (44,6\,\%) donde \texttt{gemini} apenas marca (3,8\,\%),
mientras que \texttt{gemini} domina en \texttt{masc\_generico}. Cada modelo tiene su
propio umbral por variable.

\begin{table}[H]
    \centering
    \begingroup
    \scriptsize
    \renewcommand{\texttt}[1]{{\ttfamily #1}}
    \ttabbox[\FBwidth]{
        \caption{Prevalencia de \textit{Sí} (\%) de cada modelo por variable, en el nivel B1 ($N=1.313$)}
        \label{tab:iris-prev-modelo-var}
    }{
        \begin{tabularx}{\textwidth}{
            >{\hsize=1.6\hsize\raggedright\arraybackslash\ttfamily}X
            >{\hsize=0.85\hsize\centering\arraybackslash}X
            >{\hsize=0.85\hsize\centering\arraybackslash}X
            >{\hsize=0.85\hsize\centering\arraybackslash}X
            >{\hsize=0.85\hsize\centering\arraybackslash}X
        }
            \toprule
            \normalfont\textbf{Variable} & \textbf{Gemini} & \textbf{GPT-4o-mini} & \textbf{GPT-5.4-nano} & \textbf{Gemma} \\
            \midrule
            lenguaje\_sexista & 3,8 & 44,6 & 0,1 & 1,1 \\
            masc\_generico & \textbf{66,9} & 19,1 & 9,7 & 36,9 \\
            sexismo\_discurso & 2,0 & 10,7 & 11,2 & 2,8 \\
            asimetria\_mujer\_hombre & 3,9 & 14,4 & 0,7 & 6,9 \\
            denominacion\_sexualizada & 0,9 & 1,2 & 0,1 & 1,4 \\
            \bottomrule
        \end{tabularx}
    }
    \endgroup
\end{table}

Un rasgo atraviesa a las cinco variables: el recall es bajo casi en todas, es decir,
los modelos detectan solo una fracción del sexismo que marcan las expertas. Esta
\textbf{infradetección generalizada} es el patrón dominante del capítulo, y se
examina de frente en la Sección~\ref{subsec:iris_errores} para dirimir si es una
limitación de los modelos o de la propia tarea.

\subsection{El modelo local: rendimiento y coste}
\label{subsec:iris_local}\label{subsec:iris_coste}

Entre los cuatro modelos, el local \texttt{gemma4:e4b} merece una lectura aparte. En
el nivel de control (B0) obtiene el \textbf{mayor $\kappa$ de los cuatro} (0,061), por
encima de las tres API de pago (0,026, 0,033 y 0,052), y ante las \textit{skills} se
comporta como los modelos de OpenAI: su acuerdo baja de 0,061 a 0,043
(Tabla~\ref{tab:iris-descomposicion}). Que un modelo abierto, gratuito y desplegable
en infraestructura propia iguale, y en $\kappa$ supere, a los servicios
comerciales tiene consecuencias prácticas para un proyecto que maneja material
sujeto a privacidad. Ahora bien, esa ventaja en acuerdo y en ausencia de tarifas
convive con un coste de otro tipo: el del cómputo.

La Tabla~\ref{tab:iris-b0b1} recoge el coste real de inferencia. La arquitectura de
agentes multiplica por 2,2 el coste en API respecto al \textit{baseline} (30,91 USD
frente a 13,83 USD para los tres modelos de pago sobre 1.313 noticias): el agente
encadena en promedio entre dos y tres llamadas por variable (una por cada acción del
bucle más el veredicto final) frente a la única llamada del control. La diferencia no
es mayor porque las llamadas de B0 son más largas, al inyectar la metodología
completa en el \textit{prompt}.

En el modelo local no hay tarifas, pero el coste reaparece como tiempo de cómputo. El
nivel B1 de \texttt{gemma4:e4b} se ejecutó sobre la granja de GPU (NVIDIA RTX 3090 y
4090) del Departamento de Teoría de la Señal y Comunicaciones (TSC) de la Universidad
Carlos III de Madrid, con un rendimiento medido de unos 44 segundos por variable,
esto es, en torno a 220 segundos por artículo completo. Repartiendo la carga en
cuatro GPU, cada corrida sobre el corpus de evaluación tardó alrededor de veinte
horas. En una sola GPU, el mismo cómputo superaría las ochenta horas. Así, el modelo
local es competitivo en acuerdo y prescinde de tarifas, pero su despliegue agéntico
exige infraestructura dedicada: el coste no desaparece, cambia de dinero a cómputo.
La Figura~\ref{fig:iris-acuerdo-dolar} resume ese compromiso: situado el acuerdo del
sistema desplegado (B1) frente a su coste monetario, el modelo local ocupa la esquina
más favorable (máximo acuerdo por dólar), con la salvedad de que su coste se traslada
al cómputo. Las implicaciones prácticas para IRIS se retoman en la discusión
(Sección~\ref{sec:discussion}).

\begin{figure}[H]
    \centering
    \includegraphics[width=\textwidth]{fig_acuerdo_por_dolar_blanco.pdf}
    \caption{Acuerdo (coeficiente $\kappa$ en B1) frente al coste de inferencia sobre
    el corpus completo, por modelo ($N=1.313$). El modelo local ofrece el mejor
    acuerdo por dólar (coste monetario nulo). Su coste se traslada al tiempo de
    cómputo en infraestructura propia.}
    \label{fig:iris-acuerdo-dolar}
\end{figure}

\subsection{¿Un límite del modelo o de la tarea?}
\label{subsec:iris_errores}

El bajo acuerdo absoluto plantea la pregunta que condiciona la interpretación de
todo el capítulo: ¿refleja una limitación de los modelos, o un límite intrínseco de
la propia tarea? Tres análisis complementarios apuntan en la misma dirección.

\subsubsection{Heterogeneidad entre equipos de anotación}

El corpus registra, para cada pieza, la persona que la codificó. Las 1.313
noticias anotadas se reparten entre dos equipos: Indexa (548 piezas) y UCM3
(765). Aunque la ausencia de doble codificación
(Sección~\ref{subsec:limitacion_fiabilidad}) impide calcular la fiabilidad
intercodificadora, sí es posible medir el acuerdo del sistema con cada equipo
por separado, a modo de análisis de sensibilidad. La
Tabla~\ref{tab:iris-sensibilidad} recoge el resultado para el nivel B1.

\begin{table}[H]
    \centering
    \begingroup
    \scriptsize
    \renewcommand{\texttt}[1]{{\ttfamily #1}}
    \ttabbox[\FBwidth]{
        \caption{Análisis de sensibilidad: acuerdo (B1) según el equipo de anotación}
        \label{tab:iris-sensibilidad}
    }{
        \begin{tabularx}{\textwidth}{
            >{\hsize=1.5\hsize\raggedright\arraybackslash}X
            >{\hsize=1.0\hsize\raggedright\arraybackslash}X
            >{\hsize=0.85\hsize\centering\arraybackslash}X
            >{\hsize=0.85\hsize\centering\arraybackslash}X
            >{\hsize=0.8\hsize\columncolor{gray!15}\centering\arraybackslash}X
        }
            \toprule
            \textbf{Modelo} & \textbf{Subconjunto} & \textbf{Exactitud} & \textbf{F1 macro} & \boldmath$\kappa$ \\
            \midrule
            \multirow{3}{*}{\texttt{gemini-3.1-flash-lite}}
              & Todas (1.313)   & 0,661 & 0,443 & 0,066 \\
              & Solo Indexa (548) & \textbf{0,741} & \textbf{0,503} & \textbf{0,118} \\
              & Solo UCM3 (765)  & 0,604 & 0,398 & 0,043 \\
            \midrule
            \multirow{3}{*}{\texttt{gpt-4o-mini}}
              & Todas           & 0,612 & 0,438 & 0,014 \\
              & Solo Indexa     & 0,693 & 0,482 & 0,042 \\
              & Solo UCM3       & 0,554 & 0,398 & 0,001 \\
            \midrule
            \multirow{3}{*}{\texttt{gpt-5.4-nano}}
              & Todas           & 0,569 & 0,364 & 0,004 \\
              & Solo Indexa     & 0,661 & 0,406 & 0,014 \\
              & Solo UCM3       & 0,502 & 0,323 & $-0,001$ \\
            \midrule
            \multirow{3}{*}{\texttt{gemma4:e4b}}
              & Todas           & 0,605 & 0,411 & 0,043 \\
              & Solo Indexa     & 0,695 & 0,461 & 0,068 \\
              & Solo UCM3       & 0,541 & 0,369 & 0,032 \\
            \bottomrule
        \end{tabularx}
    }
    \endgroup
\end{table}

El acuerdo con Indexa es sistemáticamente muy superior al obtenido con UCM3, en
los cuatro modelos y en las tres métricas de forma coherente. Para
\texttt{gemini-3.1-flash-lite}, restringir la evaluación a Indexa prácticamente
duplica el coeficiente $\kappa$ (de 0,066 a 0,118) y eleva la exactitud
ocho puntos. Contra UCM3 en solitario, $\kappa$ desciende hacia el azar (hasta
valores negativos en \texttt{gpt-5.4-nano}). \textbf{La magnitud de esta
variabilidad entre equipos es comparable a la diferencia entre modelos}, lo que
indica que el origen de la anotación pesa en el acuerdo tanto como la elección
del modelo.

El origen de esa diferencia está en la anotación, no en el modelo. La
Tabla~\ref{tab:iris-equipos-var} muestra, por variable, la proporción de \textit{Sí}
que marca cada equipo, y deja ver que la misma variable se codifica de forma muy
distinta. El caso extremo es \texttt{sexismo\_discurso}: Indexa marca el fenómeno en
el 11,5\,\% de las piezas y UCM3 en el 65,2\,\%, casi seis veces más. No es ruido del
modelo, son dos criterios humanos difícilmente conciliables sobre una misma variable.
Como esa variabilidad entre equipos iguala o supera la que separa a los modelos,
constituye evidencia directa del techo de la tarea, sin necesitar doble codificación.

\begin{table}[H]
    \centering
    \begingroup
    \scriptsize
    \renewcommand{\texttt}[1]{{\ttfamily #1}}
    \ttabbox[\FBwidth]{
        \caption{Prevalencia de \textit{Sí} y acuerdo (\texttt{gemini} B1) por equipo y variable. La misma variable se codifica de forma muy distinta según el equipo.}
        \label{tab:iris-equipos-var}
    }{
        \begin{tabularx}{\textwidth}{
            >{\hsize=1.7\hsize\raggedright\arraybackslash\ttfamily}X
            >{\hsize=0.9\hsize\centering\arraybackslash}X
            >{\hsize=0.9\hsize\centering\arraybackslash}X
            >{\hsize=0.9\hsize\centering\arraybackslash}X
            >{\hsize=0.8\hsize\centering\arraybackslash}X
            >{\hsize=0.8\hsize\centering\arraybackslash}X
        }
            \toprule
             & \multicolumn{2}{c}{\normalfont\% \textit{Sí}} & & \multicolumn{2}{c}{\normalfont$\kappa$} \\
            \cmidrule(lr){2-3}\cmidrule(lr){5-6}
            \normalfont\textbf{Variable} & \textbf{Indexa} & \textbf{UCM3} & \textbf{Ratio} & \textbf{Indexa} & \textbf{UCM3} \\
            \midrule
            lenguaje\_sexista & 74,6 & 85,5 & 1,1$\times$ & 0,026 & 0,009 \\
            masc\_generico & 71,4 & 88,0 & 1,2$\times$ & 0,277 & 0,223 \\
            sexismo\_discurso & \textbf{11,5} & \textbf{65,2} & \textbf{5,7$\times$} & 0,065 & $-0,003$ \\
            asimetria\_mujer\_hombre & 3,3 & 8,2 & 2,5$\times$ & 0,202 & $-0,027$ \\
            denominacion\_sexualizada & 6,6 & 12,4 & 1,9$\times$ & 0,019 & 0,013 \\
            \bottomrule
        \end{tabularx}
    }
    \endgroup
\end{table}

Debe subrayarse que este resultado \emph{no} autoriza a reportar únicamente el
subconjunto más favorable: al no existir doble codificación, no puede
determinarse qué equipo refleja mejor el constructo, y la diferencia está además
confundida con el hecho de que cada equipo codificó piezas distintas. El
análisis acota el techo alcanzable, pero no lo atribuye a un error de anotación.

\subsubsection{Consenso de los modelos frente al \textit{ground truth}}

Un segundo análisis examina las celdas en las que \textbf{los cuatro modelos
coinciden entre sí} pero difieren de la anotación experta: son las más
diagnósticas, pues descartan un error achacable a un solo modelo. De
las 6.565 celdas evaluadas (1.313 piezas $\times$ 5 variables), 1.364 presentan
consenso de los cuatro modelos en contra del \textit{ground truth}. Su dirección
es prácticamente unánime: en \textbf{1.363 casos} los cuatro modelos responden \textit{No}
donde la anotación dice \textit{Sí}, y solo en 1 caso ocurre lo contrario. La
infradetección observada por variable es, por tanto, un \textbf{sesgo sistemático},
no aleatorio: los modelos detectan de forma consistente menos sexismo que las
anotadoras.


\subsubsection{El sistema como una anotadora más}
\label{subsec:iris_equipo_ia}

Los análisis anteriores comparan cada modelo contra la anotación experta tomada
como referencia. Un enfoque complementario invierte la pregunta: en vez de medir
cuánto se equivoca el modelo respecto a un patrón fijo, lo sitúa como \textbf{una
anotadora más} dentro del espacio de criterios de las personas que codificaron el
corpus. La cuestión deja de ser ``¿acierta el modelo?'' y pasa a ser ``¿dónde cae su
criterio respecto al de las expertas?''. El sistema se trata así como un tercer
equipo de anotación, junto a Indexa y UCM3.

Como cada noticia la codificó una sola persona, no hay solape humano-humano. Lo que
sí compara a todos sin solape es la \textbf{prevalencia de \textit{Sí}}: la proporción de
piezas en que cada codificador, humano o modelo, marca la variable. La
Tabla~\ref{tab:iris-criterios} ordena por esa proporción, promediada sobre las cinco
variables, a los ocho anotadores humanos y a los cuatro modelos evaluados.

\begin{table}[H]
    \centering
    \begingroup
    \scriptsize
    \renewcommand{\texttt}[1]{{\ttfamily #1}}
    \ttabbox[\FBwidth]{
        \caption{Espacio de criterios: prevalencia media de \textit{Sí} por anotador (humano y modelo)}
        \label{tab:iris-criterios}
    }{
        \begin{tabularx}{\textwidth}{
            >{\hsize=1.5\hsize\raggedright\arraybackslash}X
            >{\hsize=0.8\hsize\centering\arraybackslash}X
            >{\hsize=0.7\hsize\centering\arraybackslash}X
            >{\hsize=1.0\hsize\columncolor{gray!15}\centering\arraybackslash}X
        }
            \toprule
            \textbf{Anotador} & \textbf{Equipo} & \boldmath$n$ & \textbf{Prev. \textit{Sí} media} \\
            \midrule
            UCM3 3 & UCM3 & 363 & 0,639 \\
            UCM3 4 & UCM3 & 166 & 0,465 \\
            Indexa 1 & Indexa & 127 & 0,408 \\
            Indexa 5 & Indexa & 50 & 0,396 \\
            UCM3 5 & UCM3 & 235 & 0,371 \\
            Indexa 2 & Indexa & 128 & 0,325 \\
            Indexa 3 & Indexa & 231 & 0,289 \\
            Indexa 4 & Indexa & 12 & 0,283 \\
            \midrule[\heavyrulewidth]
            \texttt{gpt-4o-mini} & IA & 1.313 & 0,180 \\
            \texttt{gemini-3.1-flash-lite} & IA & 1.313 & 0,155 \\
            \texttt{gemma4:e4b} & IA & 1.313 & 0,099 \\
            \texttt{gpt-5.4-nano} & IA & 1.313 & 0,043 \\
            \bottomrule
        \end{tabularx}
    }
    \endgroup
\end{table}

Los cuatro modelos ocupan el extremo permisivo de la escala, por debajo de la
anotadora humana más contenida. No se sitúan entre Indexa y UCM3: caen fuera del
rango humano completo. Es la infradetección sistemática de la subsección anterior,
leída ahora sobre una escala común a personas y máquinas. Como cada persona
codificó piezas distintas, su prevalencia mezcla dos factores: el criterio propio y
la muestra que le tocó. Para separarlos, se aplicó cada modelo, como instrumento
común, a las mismas piezas que vio cada anotadora. El orden apenas cambia: UCM3-3
sigue siendo, con diferencia, la anotadora más permisiva.

Un segundo análisis mide el acuerdo \emph{entre} modelos, que sí tienen solape
total porque los cuatro codificaron las 1.313 piezas. El resultado es inesperado y
se recoge en la Tabla~\ref{tab:iris-intraia}.

\begin{table}[H]
    \centering
    \begingroup
    \scriptsize
    \renewcommand{\texttt}[1]{{\ttfamily #1}}
    \ttabbox[\FBwidth]{
        \caption{Acuerdo entre modelos ($\kappa$ media sobre las cinco variables), por nivel}
        \label{tab:iris-intraia}
    }{
        \begin{tabularx}{\textwidth}{
            >{\hsize=1.7\hsize\raggedright\arraybackslash}X
            >{\hsize=0.65\hsize\centering\arraybackslash}X
            >{\hsize=0.65\hsize\columncolor{gray!15}\centering\arraybackslash}X
        }
            \toprule
            \textbf{Par de modelos} & \boldmath$\kappa$ \textbf{(B0)} & \boldmath$\kappa$ \textbf{(B1)} \\
            \midrule
            \texttt{gemini} $\leftrightarrow$ \texttt{gpt-4o-mini}     & 0,079 & 0,052 \\
            \texttt{gemini} $\leftrightarrow$ \texttt{gpt-5.4-nano}    & 0,067 & 0,014 \\
            \texttt{gpt-4o-mini} $\leftrightarrow$ \texttt{gpt-5.4-nano} & \textbf{0,280} & 0,084 \\
            \midrule
            \texttt{gemma} $\leftrightarrow$ \texttt{gemini}          & 0,121 & \textbf{0,151} \\
            \texttt{gemma} $\leftrightarrow$ \texttt{gpt-4o-mini}     & 0,131 & 0,104 \\
            \texttt{gemma} $\leftrightarrow$ \texttt{gpt-5.4-nano}    & 0,111 & 0,025 \\
            \bottomrule
        \end{tabularx}
    }
    \endgroup
\end{table}

Los modelos concuerdan entre sí tan poco como con las personas ($\kappa$ entre
0,014 y 0,151 en B1). Un ``equipo de IA'' con criterio compartido no existe: cada
modelo tiene su propio patrón, y lo único que tienen en común es la dirección del
sesgo, no los casos concretos. Esto desmonta la lectura fácil de que ``la IA tiene
un sesgo'' uniforme. Hay cuatro sesgos distintos que solo coinciden en ser
permisivos. El modelo local \texttt{gemma4:e4b} es el que más solape de criterio
mantiene con el resto (sus pares llegan a 0,151 con \texttt{gemini}), pero ni
siquiera ese máximo alcanza un acuerdo apreciable. Repetir el análisis sin la
arquitectura de agentes (nivel B0) deja el patrón intacto, de modo que el sesgo es
del modelo y no del sistema de \textit{skills}. Un matiz sí cambia: las
\textit{Agent Skills} \textbf{reducen} el acuerdo entre modelos en cinco de los seis
pares (el par \texttt{gpt-4o-mini} y \texttt{gpt-5.4-nano} baja de 0,280 a 0,084),
lo que reitera que la arquitectura acentúa las diferencias de cada modelo en lugar
de acercarlos a un criterio común.


Los tres análisis convergen. La baja concordancia \textbf{no debe leerse como
incapacidad de los modelos para detectar sexismo}, sino como el producto de un
criterio experto amplio y en parte en disputa, de una heterogeneidad sustancial
entre anotadores (del mismo orden que la distancia entre humanas y máquinas, con
prevalencias medias que van de 0,283 a 0,639) y de la ambigüedad de algunas
definiciones del codebook. El coeficiente $\kappa$ mide, en buena medida, esa
distancia de criterio: \textbf{el techo lo pone la tarea, no el modelo}.

